\section{Theoretical Results}
\label{sec:theory}

%This section presents theoretical guarantees. Detailed proofs are given in Appendix \ref{app:theory-proofs}.

%We give a theoretical account of AR-MAD as an agent-permutation equivariant routing rule. The guiding principle is that agent labels are arbitrary: if the agents are relabeled, the communication pattern and intermediate states should be relabeled in the same way, while the distribution of the final answer remains unchanged. This parallels permutation-equivariant mechanism design, where symmetrization preserves the objective while reducing unnecessary dependence on arbitrary labels~\citep{qin2022benefits}.

\paragraph{Formalization.}

Let $\mathfrak S_n$ be the symmetric group over agents. For
$\sigma\in\mathfrak S_n$, write $\sigma s$ for the debate state obtained by
renaming each agent $i$ as $\sigma(i)$, and write
\(
    \sigma E = \{(\sigma u,\sigma v):(u,v)\in E\}
\)
for the corresponding relabeling of a directed edge set. We assume the candidate
family is closed under relabeling:
\(
    E\in\mathcal C(s)
    \Longleftrightarrow
    \sigma E\in\mathcal C(\sigma s)
\).
This holds for the full role-assignment family induced by a fixed base graph,
since each candidate is obtained by assigning concrete agents to the same set of
roles. Let $m=|E_0|=nk$ denote the number of critique edges in every candidate
graph.

\paragraph{Equivariance.}
PEAR scores each candidate edge set by
\(
    \alpha_T\widetilde T(E\mid s)
    -
    \alpha_I\widetilde I(E\mid s)
    -
    \alpha_L\widetilde L(E\mid s)\),
where $\widetilde T$ rewards confident disagreement toward uncertain targets,
$\widetilde I$ penalizes allocating out-degree to previously influential agents,
and $\widetilde L$ penalizes low-confidence sources. We thus have the following theorem that proves PEAR is an agent-equivariant router. The proof is in Appendix \ref{app:armad-equivariant}. %The softmax router is
%\[
%    Q_\alpha(E\mid s)
%    =
%    \frac{\exp(S_\alpha(E\mid s)/\tau)}
%    {\sum_{E'\in\mathcal C(s)}\exp(S_\alpha(E'\mid s)/\tau)} .
%\]

\begin{theorem}[Agent-permutation equivariance]
\label{thm:armad-equivariant}
Assume the candidate family is closed under agent relabeling. Then, for x
$\sigma\in\mathfrak S_n$,
%\begin{gather*}    
    $S_\alpha(\sigma E\mid \sigma s)=S_\alpha(E\mid s)$, %\\
    $Q_\alpha(\sigma E\mid \sigma s)=Q_\alpha(E\mid s)$.
%\end{gather*}
\end{theorem}

%The theorem proves that AR-MAD is an agent-equivariant router. %In the deterministic limit $\tau\to0$, the same statement holds for argmax routing whenever ties are broken equivariantly, for example uniformly among maximizers.

%formalizes the role of the base graph: it fixes a sparse communication budget, while the role-to-agent assignment searches over permutations of this budget without privileging arbitrary agent labels.

\paragraph{Symmetrized routing.}
For comparison, consider any possibly non-equivariant routing policy over full
routing transcripts. Let
\(
    \Gamma=(E^{(1)},\ldots,E^{(R)})
\)
be a sequence of selected edge sets, and let $q(\Gamma\mid s^{(0)})$ be the
distribution over routing transcripts from initial state $s^{(0)}$. Define
orbit-averaged router
\(
    (\mathcal P q)(\Gamma\mid s^{(0)})
    =
    \frac{1}{n!}
    \sum_{\sigma\in\mathfrak S_n}
    q(\sigma\Gamma\mid \sigma s^{(0)})\),
where
$\sigma\Gamma=(\sigma E^{(1)},\ldots,\sigma E^{(R)})$.
Let $A(\Gamma,s^{(0)})\in[0,1]$ denote the final-answer accuracy induced by
routing $\Gamma$ from initial state $s^{(0)}$.

\begin{theorem}[Accuracy preservation]
\label{thm:accuracy-preservation}
Suppose the distribution of initial debate states is exchangeable, the agent
update kernels are permutation-equivariant, and the final aggregation rule is
permutation-invariant. Then
\[
    \mathbb E_{s^{(0)},\Gamma\sim \mathcal Pq}
    \left[A(\Gamma,s^{(0)})\right]
    =
    \mathbb E_{s^{(0)},\Gamma\sim q}
    \left[A(\Gamma,s^{(0)})\right].
\]
Moreover, $\mathcal Pq$ is agent-equivariant:
\[
    (\mathcal Pq)(\sigma\Gamma\mid \sigma s^{(0)})
    =
    (\mathcal Pq)(\Gamma\mid s^{(0)})\text{, }
    \forall \sigma\in\mathfrak S_n .
\]
\end{theorem}

Thus, under exchangeable agent populations, enforcing equivariance does not
change expected final accuracy; it only removes arbitrary dependence on the agent
names. The proof is in Appendix \ref{app:accuracy-preservation}.

\paragraph{Generalization.}
We next formalize the complexity benefit of equivariant routers. Let
$\mathcal Q$ be a class of one-round routers $q(\cdot\mid s)$, and define
\(
    d(q,q')
    =
    \sup_s
    \left\|
    q(\cdot\mid s)-q'(\cdot\mid s)
    \right\|_1 
\).
Let
\(
    \mathcal P\mathcal Q=\{\mathcal Pq:q\in\mathcal Q\}
\)
be the orbit-projected class. We use covering number to measure the hypothesis complexity; see Definition \ref{def:covering-number} in Appendix \ref{app:complexity}.



\begin{lemma}[Covering number control]
\label{lem:covering-projection}
For every $\epsilon>0$, the covering number satisifies,
\(
    \mathcal N(\mathcal P\mathcal Q,\epsilon,d)
    \le
    \mathcal N(\mathcal Q,\epsilon,d)
\).
\end{lemma}

\begin{theorem}[Generalization bound]
\label{thm:generalization}
Let $z_1,\ldots,z_{N_{\rm ex}}$ be i.i.d. task instances, and let
$F(q;z)\in[0,1]$ be an evaluation functional, such as final-answer accuracy,
one-round correction, or negative routing regret. Suppose $F$ is
$L$-Lipschitz in $q$ under $d$:
\(
    |F(q;z)-F(q';z)|\le Ld(q,q')\),
\(    \forall q,q',z 
\).
Define
\(
    \mathcal L(q)=\mathbb E_z[F(q;z)]\) and
\(    \widehat{\mathcal L}_{N_{\rm ex}}(q)
    =
    \frac{1}{N_{\rm ex}}
    \sum_{\ell=1}^{N_{\rm ex}}F(q;z_\ell)\).
Then, for any $\epsilon>0$ and $\delta\in(0,1)$, with probability at least
$1-\delta$,
%\[
%\sup_{q\in\mathcal P\mathcal Q}
%\left|
%\mathcal L(q)-\widehat{\mathcal L}_{N_{\rm ex}}(q)
%\right|
%\le
%2L\epsilon
%+
%\sqrt{
%\frac{
%\log\left(2\mathcal N(\mathcal P\mathcal Q,\epsilon,d)/\delta\right)
%}{
%2N_{\rm ex}
%}}
%.
%\]
%Consequently, by Lemma~\ref{lem:covering-projection},
\begin{align*}    
%&\sup_{q\in\mathcal P\mathcal Q}
&\left|
\mathcal L(q)-\widehat{\mathcal L}_{N_{\rm ex}}(q)
\right|\\
\le&
2L\epsilon
+
\sqrt{
\frac{
\log\left(2\mathcal N(\mathcal Q,\epsilon,d)/\delta\right)
}{
2N_{\rm ex}
}}
.
\end{align*}
\end{theorem}

Lemma \ref{lem:covering-projection} and Theorem~\ref{thm:generalization} show that restricting to equivariant routers cannot increase the covering complexity of the routing class. PEAR therefore combines a sparse communication budget with an equivariant search space, reducing both per-round debate cost and the effective complexity of the router. Proofs are in Appendix \ref{app:complexity}.

\paragraph{One-round correction and influence-balancing.} More theoretical results on one-round correction and influence-balancing are given in Appendix \ref{app:deferred-theory}.